
\documentclass[english]{article}
%\documentclass[aps,pra,twocolumn,english]{revtex4}
%\documentclass[english]{iopart}
\usepackage[T1]{fontenc}
\usepackage[latin9]{inputenc}
\usepackage{babel}
\usepackage{graphicx,amsmath,amssymb,color}
\usepackage[normalem]{ulem}
\usepackage{amsfonts}
\usepackage[toc,page]{appendix}
%\usepackage[ocgcolorlinks,colorlinks=true,linkcolor=blue,citecolor=red]{hyperref}
\usepackage{hyperref}
\usepackage{latexsym}
\usepackage{amsfonts}
\usepackage{algpseudocode}
\usepackage{amsthm}
\usepackage{mathrsfs}
\usepackage{natbib}
\usepackage{color,verbatim}
\usepackage{psfrag}
\bibliographystyle{unsrt}
%\bibliographystyle{acm}
%\bibliographystyle{unsrtnat}
%\usepackage[style=numeric]{biblatex}

\usepackage{subcaption} % allows for subfigures


\usepackage[svgnames]{xcolor}
\usepackage{tikz}

\usepackage{algorithm}


%\input{tikzit/tikz_preamble.tex}  % load separately defined tikz preamble


\newcommand{\bra}[1]{\langle#1 |}
\newcommand{\ket}[1]{|#1 \rangle}
\newcommand{\braket}[2]{\left \langle #1 \mid #2 \right \rangle}
\newcommand{\bigbraket}[2]{\left \langle #1 \middle\vert #2 \right \rangle}
\newcommand{\ketbra}[2]{\vert #1 \rangle \! \langle #2 \vert}
\newcommand{\overlap}[2]{\left \langle #1 | #2 \right\rangle}
\newcommand{\average}[1]{\langle #1 \rangle}
\newcommand{\sandwich}[3]{\left \langle #1 \mid #2 \mid #3 \right\rangle}
\newcommand{\innerprod}[2]{\left \langle #1 | #2 \right\rangle}

\begin{document}


\title{Non-signalling, Fourier transforms, and particle statistics}
\author{Nicholas Chancellor}


\maketitle


\today % added for drafting, remove in final version


{\small The aim of this worksheet is to have a mixture of hard and easy questions, do not feel like they need to be answered in order, feel free to skip around and come up with partial answers wherever possible; feel free to use maple/matlab/python/wolfram alpha to aid in calculations.}

\begin{enumerate}
\item Consider a quantum system which can initially be described by a density matrix $\hat{\rho}$, and has subsystems $A$ and $B$, which may or may not be entangled. [Hint: it may be worth recalling some of the properties of the trace operation, in particular the fact that it is cyclic $\mathrm{Tr}[PQR]=\mathrm{Tr}[QRP]\quad \forall P,R,Q$]
\begin{itemize}
\item Show that the measured value of an observable $\hat{O}_A=\hat{O}\otimes 1_B$ on $A$ is unaffected by applying a unitary to system $B$ ($\hat{U}_B=1_A\otimes \hat{U}$).
\item Now consider that we apply a projective measurement on $B$, which projects the system into a single state (represented mathematically by applying $\hat{\Pi}_B=1_A \otimes \ketbra{\phi}{\phi}$), will this change our observed values? Is this compatible with relativity?
\item What if we instead consider a projective measurement with an unknown outcome each potential outcome performing a projection $\hat{\Pi}_{B,j}$? [hint: remember that quantum mechanics is a linear theory and density matrices obey the superposition principle, the expectation of a sum of density matrices is the sum of their individual expectations. Furthermore for probability to be conserved we must have $\sum_j \hat{\Pi}_{B,j}=1_{A}\otimes 1_{B}$ ] What does this imply for the ability to send information instantaneously by taking measurements if the outcome is not communicated?
\item Now consider the completely general case, where system $B$ is operated on by a general open system process described by the Lindblad equation,
\begin{equation}
\frac{\partial \hat{\rho}}{\partial t}=-\frac{i}{\hbar}\left[\hat{\rho}, \hat{H}\right]+\sum_j\gamma_j\left(\hat{L}_j\hat{\rho}\hat{L}_j^\dagger-\frac{1}{2}\hat{L}_j^\dagger \hat{L}_j\hat{\rho}-\frac{1}{2}\hat{\rho}\hat{L}_j^\dagger \hat{L}_j\right)
\end{equation}
 what property of $\hat{H}$, $\hat{L}_j$, and $\gamma_j\ge 0$  is sufficient to guarantee that the measurement value is unchanged? How could this be ensured if we are allowed completely a general $\hat{O}$ in $\hat{O}_A=\hat{O}\otimes 1_B$?  what does this mean physically?
\end{itemize}
\item For this question we define the Fourier transform as $\tilde{g}(\omega)=\mathcal{F}[g(t)](\omega)= \int_{-\infty}^\infty dt\, g(t)\exp(-i\omega t)$ and the inverse Fourier transform as $g(t)=\mathcal{F}^{-1}[\tilde{g}(\omega)](t)=\frac{1}{2\pi} \int_{-\infty}^\infty d\omega\, \tilde{g}(\omega)\exp(i\omega t)$ there are several widely used conventions/notations here but we pick this one for concreteness
\begin{itemize}
\item Show that $\mathcal{F}[g(a t)](\omega)=\mathcal{F}[g(t)](\frac{\omega}{a})$, and that it follows that the products of the second moment will be invariant under this scaling
\begin{align}
\left(\int_{-\infty}^{\infty}dt\, t^2 g(a t)\right)\times \left(\int_{-\infty}^{\infty}d\omega\,\omega^2 \mathcal{F}[g(a t)](\omega) \right)=\\
\left(\int_{-\infty}^{\infty}dt\, t^2 g(t)\right)\times \left(\int_{-\infty}^{\infty}d\omega\,\omega^2 \mathcal{F}[g(t)](\omega) \right)
\end{align}
\item Prove the convolution theorem, in other words show that 
\begin{equation}
\mathcal{F}[g(t)h(t)](\omega)=\int_{-\infty}^\infty d\omega' \mathcal{F}[g(t')](\omega')\mathcal{F}[h(t)](\omega-\omega'),
\end{equation}
 and argue the equivalent holds for the inverse Fourier transform. [hint: the following identity may be helpful $ \int_{-\infty}^\infty d\omega'  \exp(i\omega'(t-t'))=\delta(t-t')$, where $\delta$ is the Dirac delta]
\item By the convolution theorem argue that the frequency distribution of a wavepacket of shape $g(t)$ put through a shutter which opens at time $-t_0$ and closes at time $t_0$ will be $\int_{-\infty}^\infty d\omega'\mathcal{F}[g(t')](\omega') \frac{2\sin((\omega-\omega') t_0)}{\omega-\omega'}$
\item Use the previous result to argue that the frequency distribution attained by applying a shutter to a single mode at $\omega_p$ will give a frequency distribution of $\tilde{g}(\omega)=\frac{2\sin((\omega-\omega_p) t_0)}{\omega-\omega_p}$
\item Find the Fourier transform of the following function, corresponding to a system which is turned on and has an exponential decay envelope defined by $\gamma$, where $\Theta$ is the Heviside Theta function)
\begin{equation}
\Theta(t)\exp(i \omega_p t-\gamma t) 
\end{equation}
Show that it gives
\begin{equation}
 \tilde{g}(\omega)=-\frac{1}{i(\omega_p-\omega)-\gamma}
 \end{equation}
separate into real and imaginary parts and argue that in the limit $\gamma \rightarrow 0$, the real part of this function becomes a Dirac delta centred at $\omega=\omega_p$, while the imaginary part remains a function with finite width, what does this finite width mean physically? [hint: a useful identity which can be shown by contour integration (or Wolfram alpha), that  $\int_{-\infty}^{\infty}d\omega\frac{\gamma}{\omega^2+\gamma^2}=\pi$ independent of $\gamma$]
\end{itemize}
\item Consider the Hong-Ou-Mandel setup (single beamsplitter with one particle in each input mode) but with labelled particles $\hat{a}$ and $\hat{b}$, we will label the modes $0$, $1$, $2$, $3$ as in the notes
\begin{itemize}
\item Show that if we start with $\hat{a}^\dagger_0\hat{b}^\dagger_1\ket{0}$ the probabilities of each configuration will be no different than if the directions were decided by a classical coin flip
\item Now consider a symmetric superposition of $\hat{a}$ and $\hat{b}$ particles $\frac{1}{\sqrt{2}}(\hat{a}^\dagger_0\hat{b}^\dagger_1+\hat{a}^\dagger_1\hat{b}^\dagger_0)\ket{0}$, show that this recovers the Hong-Ou-Mandel result, what does this mean physically? [hint: the creation operators for $a$ and $b$ will commute]
\item Now consider taking an anti-symmetric superposition  $\frac{1}{\sqrt{2}}(\hat{a}^\dagger_0\hat{b}^\dagger_1-\hat{a}^\dagger_1\hat{b}^\dagger_0)\ket{0}$, what kind of particle would this correspond to and what does the result mean physically?
\end{itemize}
\end{enumerate}







\end{document}